\begin{exo}{}
	On considère un triangle $EFG$ rectangle en $F$ tel que :
	\cmath{EF = \dfrac{3}{4} \times 10^2 \quad \text{et} \quad FG = \dfrac{1}{2} \times 10^3}
	\begin{center}
		\begin{tikzpicture}[scale=0.012, >=Stealth]
			\coordinate (F) at (0,0);
			\coordinate (E) at (0,75);
			\coordinate (G) at (500,0);
			\draw[very thick, black] (E) -- (F) -- (G) -- cycle;
			\draw[thick] (0,20) -- (20,20) -- (20,0);
			\node[above left] at (E) {$E$};
			\node[below left] at (F) {$F$};
			\node[below right] at (G) {$G$};
			\node[above right] at (250,37.5) {$?$};
		\end{tikzpicture}
	\end{center}
	\begin{enumerate}
		\item Déterminer l'écriture décimale des longueurs $EF$ et $FG$.
		\item En utilisant le théorème de Pythagore, calculer la valeur exacte de $EG^2$. On donnera le résultat sous la forme d'une fraction simplifiée multipliée par une puissance de $10$.
		\item En déduire la longueur exacte de l'hypoténuse $EG$ sous la forme $a\sqrt{b}$ où $a$ et $b$ sont des entiers.
	\end{enumerate}
\end{exo}
